% TeX'ing this file requires that you have AMS-LaTeX 2.0 installed
% as well as the rest of the prerequisites for REVTeX 4.0
%
% See the REVTeX 4 README file
% It also requires running BibTeX. The commands are as follows:
%
%  1)  latex apssamp.tex
%  2)  bibtex apssamp
%  3)  latex apssamp.tex
%  4)  latex apssamp.tex
%
\documentclass[prb,preprintnumbers,amsmath,amssymb, 11pt]{revtex4}
%\documentclass[preprint,showpacs,showkeys,preprintnumbers,amsmath,amssymb]{revtex4}

% Some other (several out of many) possibilities
%\documentclass[preprint,aps]{revtex4}
%\documentclass[aps, two column, amsmath,amssymb,floatfix]{revtex4}
%\documentclass[showkeys,showpacs,amsmath,amssymb,onecolumn,superscriptaddress,prl]{revtex4-1}% Physical Review B  
% \documentclass[aps,prl,reprint,showpacs,floatfix,superscriptaddress, onecolumn]{revtex4-2}

\usepackage{amsmath,amsthm,amssymb}
\usepackage{graphicx}% Include figure files
\usepackage{dcolumn}% Align table columns on decimal point
\usepackage{bm}% bold math
\usepackage{color}
\usepackage{epsfig}
\usepackage{multirow}
\usepackage{mathrsfs}
\usepackage{hyperref}
\usepackage{cleveref}
\usepackage{epstopdf}
\usepackage{subfigure}
\usepackage{autobreak}

%Macros for mathematical notations

\newcommand{\V}[1]{\boldsymbol{#1}} %# vector
\newcommand{\M}[1]{\boldsymbol{#1}} %# matrix
\newcommand{\Set}[1]{\mathbb{#1}} %# set
\newcommand{\D}[1]{\Delta#1} %# \D{t} for time step size
\renewcommand{\d}[1]{\delta#1} %# \d{t} for small increment
\newcommand{\norm}[1]{\left\Vert #1\right\Vert } % norm
\newcommand{\abs}[1]{\left|#1\right|} %abs

\newcommand{\grad}{\M{\nabla}} %gradient
\newcommand{\av}[1]{\left\langle #1\right\rangle } %take average

\newcommand{\sM}[1]{\M{\mathcal{#1}}} %matrix in mathcal font
\newcommand{\dprime}{\prime\prime} % double prime
%\global\long\def\i{\iota}
%\renewcommand{\i}{\iota} %i for imaginary unit
%\renewcommand{\i}{\mathsf i} %i for imaginary unit
\newcommand{\follows}{\quad\Rightarrow\quad} %=>
\newcommand{\eqd}{\overset{d}{=}} %=^d
\newcommand{\spe}[1]{\mathscr{#1}}  %important quantities in mathscr font
\newcommand{\eps}{\epsilon}
\newcommand{\bra}[1]{\left\langle #1 \right|}
\newcommand{\ket}[1]{\left| #1 \right\rangle}
\newcommand{\braket}[2]{\left\langle #1 \middle| #2 \right\rangle}
\newcommand{\bracket}[3]{\left\langle #1 \middle| #2 \middle| #3 \right\rangle}



\begin{document}
\preprint{Preprint}

\title{Review of the Fock Matrix}
\author{Xuanzhao Gao}
% \date{}

\maketitle

Under the Born-Oppenheimer approximation, the Hamiltonian of the system with $N$ electrons is given by:
\begin{equation}
    \hat{H} = \sum_{i=1}^N \left( -\frac{1}{2} \nabla_{i}^2 - \sum_{A} \frac{Z_A}{\abs{\mathbf{r}_i - \mathbf{R}_A}} \right) + \sum_{i < j} \frac{1}{|\mathbf{r}_i - \mathbf{r}_j|},
\end{equation}
where~$Z_A$ is the nuclear charge of the $A$-th nucleus,~$\mathbf{R}_A$ is the position of the $A$-th nucleus.
Then the Schrödinger equation is given by:
\begin{equation}
    \hat{H} \Psi(\mathbf{r}_1, \mathbf{r}_2, \cdots, \mathbf{r}_N) = E \Psi(\mathbf{r}_1, \mathbf{r}_2, \cdots, \mathbf{r}_N),
\end{equation}
where~$E$ is the energy of the system.
To solve this equation, we assume that the wave function can be represented by a \textit{Slater determinant}~\cite{cramer2013essentials}:
\begin{equation}
    \Psi(\mathbf{r}_1, \mathbf{r}_2, \cdots, \mathbf{r}_N) = \frac{1}{\sqrt{N!}} \left| \begin{array}{cccc}
        \psi_1(\mathbf{r}_1) \alpha(r_1) & \psi_2(\mathbf{r}_1) \alpha(r_1) & \cdots & \psi_N(\mathbf{r}_1) \alpha(r_1) \\
        \psi_1(\mathbf{r}_2) \alpha(r_2) & \psi_2(\mathbf{r}_2) \alpha(r_2) & \cdots & \psi_N(\mathbf{r}_2) \alpha(r_2) \\
        \vdots & \vdots & \ddots & \vdots \\
        \psi_1(\mathbf{r}_N) \alpha(r_N) & \psi_2(\mathbf{r}_N) \alpha(r_N) & \cdots & \psi_N(\mathbf{r}_N) \alpha(r_N)
    \end{array} \right|,
\end{equation}
where~$\psi_i(\mathbf{r})$ are normalized one-electron eigenfunctions, denoted as molecular orbitals (MO),~$\mathbf{r}_i$ is the position of the $i$-th electron,~$\alpha(r_i)$ is the spin function.
Notice that the Slater determinant is antisymmetric, i.e., if two electrons are exchanged, the wave function changes its sign.
We then have the Hatree-Fock equation:
\begin{equation}\label{eq:hatree-fock}
    \hat{F}(\V{r}_1) \psi_i(\V{r}_1) = \epsilon_i \psi_i(\V{r}_1),
\end{equation}
where~$\hat{F}$ is the Fock operator,~$\epsilon_i$ is the orbital energy of the $i$-th MO.
The Fock operator is given by:
\begin{equation}
    \hat{F}(\V{r}_1) = - \frac{1}{2} \nabla^2 - \sum_{A} \frac{Z_A}{\abs{\V{r}_1 - \mathbf{R}_A}} + \sum_{j} \left[ 2 \hat{J}_j(\V{r}_1) - \hat{K}_j(\V{r}_1) \right],
\end{equation}
and 
The Coulomb operator is given by:
\begin{equation}
    \hat{J}_j(\V{r}_1) \psi_i(\V{r}_1) = \psi_i(\V{r}_1) \int \frac{\psi_j^*(\V{r}_2) \psi_j(\V{r}_2)}{|\V{r}_1 - \V{r}_2|} \mathrm{d}\V{r}_2,
\end{equation}
The exchange operator is given by:
\begin{equation}
    \hat{K}_j(\V{r}_1) \psi_i(\V{r}_1) = \psi_j(\V{r}_1) \int \frac{\psi_j^*(\V{r}_2) \psi_i(\V{r}_2)}{|\V{r}_1 - \V{r}_2|} \mathrm{d}\V{r}_2.
\end{equation}

In actual calculation, the molecular orbitals~$\psi_i(\V{r})$ are approximated as a linear combination of atomic orbitals (basis functions):
\begin{equation}
    \psi_i(\mathbf{r}) = \sum_{\nu=1}^K C_{\nu i} \varphi_{\nu}(\mathbf{r})
\end{equation}
where~$\varphi_{\nu}$ is the atomic orbital,~$K$ is the number of atomic orbitals applied,~$C_{\nu i}$ is the coefficient of the linear combination.
We can also regard~$\varphi$ as a basis set of the wave function.
There can be several different choices of the basis set, such as the \textit{Slater-type Orbital} \cite{ira2000levine}:
\begin{equation}
    \varphi(\mathbf{r}) = r^{n - 1} e^{- \zeta r} Y_{l}^m(\theta,\phi),
\end{equation}
where~$n$ is the principal quantum number,~$l$ is the azimuthal quantum number,~$m$ is the magnetic quantum number; or the \textit{Gaussian-type Orbital}:
\begin{equation}
    \varphi(\mathbf{r}) = x^a y^b z^c e^{- \alpha r^2}.
\end{equation}
The Gaussian-type orbital is more preferred, since it is easier to calculate and store.
Inserting the basis functions into the Hatree-Fock equation \eqref{eq:hatree-fock}, we have:
\begin{equation}
    \sum_{\nu=1}^K C_{\nu i} \left( \bracket{\psi_\mu}{\hat{F}}{\psi_\nu} - \epsilon_i \braket{\psi_\mu}{\psi_\nu} \right) = 0,
\end{equation}
and can be rewritten as the Roothaan equation \cite{roothaan1951new}:
\begin{equation}\label{eq:roothaan}
    \mathrm{det} \left( \mathbf{F} - \mathbf{S} \mathbf{\epsilon} \right) = 0,
\end{equation}
where~$F$ is the Fock matrix,~$S$ is the overlap matrix,~$\epsilon$ is the orbital energy.

In practice, the key is to construct the Fock matrix, given by:
\begin{equation}
    F_{\mu \nu} = h_{\mu \nu} + \sum_{\lambda \sigma} \sum_{i} 2 C_{\lambda i} C_{\sigma i} (J_{\mu \nu \lambda \sigma} - \frac{1}{2} K_{\mu \nu \lambda \sigma}),
\end{equation}
where
\begin{equation}\label{eq:core-hamiltonian}
    h_{\mu \nu} = \int \varphi_{\mu}(\mathbf{r}) \left( -\frac{1}{2} \nabla^2 - \sum_{A} \frac{Z_A}{\abs{\mathbf{r} - \mathbf{R}_A}} \right) \varphi_{\nu}(\mathbf{r}) \mathrm{d}\mathbf{r},
\end{equation}
is the core Hamiltonian, representing the kinetic energy and the interaction with the nuclei,
\begin{equation}\label{eq:coulomb-integral}
    J_{\mu \nu \lambda \sigma} = \iint \varphi_{\mu}(\mathbf{r}_1) \varphi_{\nu}(\mathbf{r}_1) \frac{1}{|\mathbf{r}_1 - \mathbf{r}_2|} \varphi_{\lambda}(\mathbf{r}_2) \varphi_{\sigma}(\mathbf{r}_2) \mathrm{d}\mathbf{r}_1 \mathrm{d}\mathbf{r}_2,
\end{equation}
is the Coulomb integral, representing the Coulomb repulsion between two electrons,
\begin{equation}\label{eq:exchange-integral}
    K_{\mu \nu \lambda \sigma} = \iint \varphi_{\mu}(\mathbf{r}_1) \varphi_{\lambda}(\mathbf{r}_1) \frac{1}{|\mathbf{r}_1 - \mathbf{r}_2|} \varphi_{\sigma}(\mathbf{r}_2) \varphi_{\nu}(\mathbf{r}_2) \mathrm{d}\mathbf{r}_1 \mathrm{d}\mathbf{r}_2,
\end{equation}
is the exchange integral, representing the exchange correlation between two electrons.
Then in actual calculations, one starts from a initial guess of the coefficients~$C_{\nu i}$, and iteratively solves the Roothaan equation Eq.~\eqref{eq:roothaan} to obtain the final coefficients.
To construct the Fock matrix, we need to evaluate the above integrals. 
Thus if there are $N$ orbitals, the computation cost will be $O(N^4)$.
Considering the base functions are rapidly decaying, this number can be greatly reduced.


\bibliography{ref}


\end{document}